% Part: incompleteness
% Chapter: arithmetization-syntax
% Section: coding-terms

\documentclass[../../../include/open-logic-section]{subfiles}

\begin{document}

\olfileid{inc}{art}{trm}
\olsection{د ترمونو کوډول}

\begin{explain}
يو ترم يوازې د سمبولونو يو ځانګړے ډول لړۍ ده: دا د ترمونو د جوړښتي
قاعدو له مخې له ثابتونو او متغيرونو څخه په استقرايي ډول جوړېږي۔ ځکه
چې د سمبولونو لړۍ د عددونو په توګه کوډېدلے شي---د سمبولونو د کوډولو
له يوې طرحې او د عددونو د لړيو د کوډولو له يوې لارې سره---نو ترمونو
ته د ګوډل شمېرو ټاکل ستونزمن نۀ دي۔ اصلي ستونزه دا ښودل دي چې يو عدد
چې د سم جوړ شوي ترم ګوډل شمېره وي، دا خاصيت ئې محاسبه کېدونکے، بلکې
بنسټيز بازګشتي دے۔
\end{explain}

!!^{variable}s او !!{constant}s تر ټولو ساده ترمونه دي، او دا ازمويل
چې $x$ د داسې يوه ترم ګوډل شمېره ده، اسانه دي: $\fn{Var}(x)$ هله
صدق کوي چې $x$ د کوم~$i$ دپاره $\Gn{\Obj v_i}$ وي۔ په بله وينا،
$x$ د اوږدوالي~$1$ يوه لړۍ ده او يوازينے غړے ئې $(x)_0$ د کوم
!!{variable}~$\Obj v_i$ کوډ دے؛ يعنې $x$ د کوم~$i$ دپاره
$\tuple{\tuple{1, i}}$ دے۔ په ورته ډول $\fn{Const}(x)$ هله صدق کوي
چې $x$ د کوم~$i$ دپاره $\Gn{\Obj c_i}$ وي۔ دا دواړه اړيکې بنسټيزې
بازګشتي دي، ځکه که داسې $i$ موجود وي، بايد $<x$ وي:
\begin{align*}
  \fn{Var}(x) & \defiff \bexists{i<x}{x = \tuple{\tuple{1, i}}}\\
  \fn{Const}(x) & \defiff \bexists{i<x}{x = \tuple{\tuple{2, i}}}
\end{align*}

\begin{prop}
\ollabel{prop:term-primrec}
هغه اړيکې $\fn{Term}(x)$ او $\fn{ClTerm}(x)$ چې په ترتيب هله صدق کوي
چې $x$ د يوه ترم يا تړلي ترم ګوډل شمېره وي، بنسټيزې بازګشتي دي۔
\end{prop}

\begin{proof}
د سمبولونو يوه لړۍ~$s$ هله ترم دے چې د ترمونو داسې لړۍ~$s_0$،
\dots، $s_{k-1}=s$ موجوده وي چې ثبتوي ترم~$s$ د ترمونو د جوړښتي
قاعدو له مخې له !!{constant}s او !!{variable}s څخه څنګه جوړ شوے دے۔
د دې د بيانولو دپاره چې داسې يوه اټکلي جوړښتي لړۍ جوړښتي قاعدې
تعقيبوي، بايد د هر $i<k$ دپاره له لاندې څخه يو حالت سم وي:
\begin{enumerate}
\item $s_i$ يو !!a{variable}~$\Obj v_j$ دے؛ يا
\item $s_i$ يو !!a{constant}~$\Obj c_j$ دے؛ يا
\item $s_i$ له هغو $n$ ترمونو $t_1$، \dots، $t_n$ څخه، چې له
  ځاے~$i$ وړاندې راغلي وي، د يوه $n$-ځايه
  !!{function}~$\Obj f^n_j$ په کارولو جوړ شوے دے۔
\end{enumerate}
د دې د ښودلو دپاره چې د ګوډل شمېرو اړونده اړيکه بنسټيزه بازګشتي
ده، دا شرط بايد په بنسټيز بازګشتي ډول بيان کړو؛ يعنې بنسټيزې
بازګشتي تابعې، اړيکې او محدود کمیت ايښودنه وکاروو۔

فرض کړئ $y$ هغه عدد دے چې لړۍ~$s_0$، \dots، $s_{k-1}$ کوډوي؛ يعنې
$y = \tuple{\Gn{s_0}, \dots, \Gn{s_{k-1}}}$۔ دا د ګوډل شمېرې~$x$
لرونکي ترم دپاره هله جوړښتي لړۍ کوډوي چې د هر $i<k$ دپاره:
\begin{enumerate}
\item $\fn{Var}((y)_i)$؛ يا
\item $\fn{Const}((y)_i)$؛ يا
\item داسې $n$ او عدد~$z=\tuple{z_1, \dots, z_n}$ شته چې هر $z_l$
  د کوم $i'<i$ دپاره له $(y)_{i'}$ سره برابر وي، او
\[
(y)_i = \Gn{\Obj f^n_j(} \concat \fn{flatten}(z) \concat \Gn{)},
\]
\end{enumerate}
او سربېره پر دې $(y)_{k-1}=x$ وي۔ (تابع $\fn{flatten}(z)$ لړۍ
$\tuple{\Gn{t_1}, \dots, \Gn{t_n}}$ په $\Gn{t_1, \dots, t_n}$
بدلوي او بنسټيزه بازګشتي ده۔)

په (3) کښې شاخصونه $j$ او $n$، د ترمونو $t_l$ ګوډل شمېرې~$z_l$،
او د لړۍ~$\tuple{z_1, \dots, z_n}$ کوډ~$z$ ټول له~$y$ څخه واړه دي۔
پورته $k$ په $\len{y}$ بدلولے شو۔ نو «$y$ د ګوډل شمېرې~$x$ لرونکي
ترم د جوړښتي لړۍ کوډ دے» په داسې ډول بيانولے شو چې وښيي دا اړيکه
بنسټيزه بازګشتي ده۔

اوس يوازې دا باور ترلاسه کول پکار دي چې پر~$y$ يو بنسټيز بازګشتي حد
شته۔ خو که $x$ د يوه ترم ګوډل شمېره وي، نو هغه بايد داسې جوړښتي لړۍ
ولري چې تر ټولو زيات $\len{x}$ ترمونه ولري (ځکه د~$s$ په جوړښتي
لړۍ کښې هر ترم بايد د~$s$ په کوم ځاے کښې پيل شي، او دوه فرعي ترمونه
په يوه ځاے نۀ شي پيلېدے)۔ البته د~$s$ د هر فرعي ترم ګوډل شمېره
$\le x$ ده۔ نو تل داسې جوړښتي لړۍ شته چې کوډ ئې
$\le p_{k-1}^{k(x+1)}$ وي، چې $k=\len{x}$ دے۔

د $\fn{ClTerm}$ دپاره يوازې د !!{variable}s فقره پرېږدئ۔
\end{proof}

\begin{prob}
وښيئ چې تابع $\fn{flatten}(z)$، چې لړۍ
$\tuple{\Gn{t_1}, \dots, \Gn{t_n}}$ په $\Gn{t_1, \dots, t_n}$
بدلوي، بنسټيزه بازګشتي ده۔
\end{prob}

\begin{prop}\ollabel{prop:num-primrec}
  تابع $\fn{num}(n) = \Gn{\num{n}}$ بنسټيزه بازګشتي ده۔
\end{prop}

\begin{proof}
  $\fn{num}(n)$ په بنسټيز بازګښت داسې تعريفوو:
  \begin{align*}
    \fn{num}(0) & = \Gn{\Obj 0}\\
    \fn{num}(n+1) & = \Gn{\prime(} \concat \fn{num}(n) \concat \Gn{)}.
  \end{align*}
\end{proof}

\end{document}
